% from ega4-17-fr.tex
\oldpage[IV-4]{71}
un corps puisqu'il est artinien, et il résulte alors de (0, 19.6.5.1) que $A$ doit être une extension finie et séparable de $k$, ce qui achève la démonstration (17.4.1).

\begin{corollary}[17.6.2]
\label{IV.17.6.2}
Soit $f:X\to Y$ un morphisme localement de présentation finie. Les conditions suivantes sont équivalentes:
\begin{enumerate}
  \item[a)] $f$ est étale.
  \item[a')] $f$ est lisse et non ramifié.
  \item[b)] $f$ est lisse et localement quasi-fini (Err$_{\mathrm{III}}$, 20).
  \item[c)] $f$ est plat et non ramifié.
  \item[c')] $f$ est plat, et toute fibre $f^{-1}(y)$ est un $k(y)$-schéma somme de spectres de corps, extensions finies et séparables de $k(y)$.
  \item[c'')] $f$ est plat, et pour tout $y\in Y$ et toute extension algébriquement close $k'$ de $k(y)$, la « fibre géométrique » $f^{-1}(y)\otimes_{k(y)}k'$ est somme de spectres de corps isomorphes à $k'$.
\end{enumerate}
\end{corollary}

Le seul point qui reste à démontrer est l'équivalence de c') et c''). Il est clair que c') entraîne c'') par changement de base (17.3.3). D'autre part, comme le morphisme projection $f^{-1}(y)\otimes_{k(y)}k'\to f^{-1}(y)$ est ouvert (2.4.10), l'hypothèse c'') entraîne que l'espace $f^{-1}(y)$ est discret, donc, pour tout point $x\in X_y=f^{-1}(y)$, l'anneau local $\sh{O}_{X_y,x}=A$ est une $k(y)$-algèbre finie, donc un anneau local artinien; en outre il résulte de c'') que $\Spec(A\otimes_{k(y)}k')$ est somme de spectres de corps isomorphes à $k'$, ce qui n'est possible que si $A$ est un corps, extension finie séparable de $k(y)$ (4.6.1).

\begin{proposition}[17.6.3]
\label{IV.17.6.3}
Soient $Y$ un préschéma localement noethérien, $f:X\to Y$ un morphisme localement de type fini, $x$ un point de $X$, $y=f(x)$. Posons $A=\sh{O}_{Y,y}$, $B=\sh{O}_{X,x}$, qui sont des anneaux locaux noethériens, et soit $k$ le corps résiduel de $A$. Alors les conditions équivalentes a) à d) de (17.6.1) sont aussi équivalentes à chacune des suivantes:
\begin{enumerate}
  \item[e)] $\widehat B$ est une $\widehat A$-algèbre formellement étale pour les topologies adiques.
  \item[e')] $\widehat B$ est un $\widehat A$-module libre et $\widehat B\otimes_{\widehat A}k$ est un corps, extension finie et séparable de $k$ (ce qui entraîne que $\widehat B$ est une $\widehat A$-algèbre finie).
\end{enumerate}
Si de plus $k(x)=k(y)$, ou si $k$ est séparablement clos, ces conditions équivalent aussi à:
\begin{enumerate}
  \item[e'')] L'homomorphisme canonique $\widehat A\to\widehat B$ est bijectif.
\end{enumerate}
\end{proposition}

L'équivalence de e) et de chacune des conditions de (17.6.1) résulte aussitôt de (17.4.4, f')) et (17.5.3, d')). Le fait que e) entraîne e') résulte de (17.4.4, f)) et de (0, 19.7.1), compte tenu de ce que $\widehat B$ est alors une $\widehat A$-algèbre finie (17.4.4) et qu'il revient donc au même de dire que $\widehat B$ est un $\widehat A$-module plat ou un $\widehat A$-module libre ($0_{\mathrm{III}}$, 10.1.3). Inversement, le fait que e') entraîne e) résulte de (17.4.4) et de (0, 19.7.1). Enfin, e') entraîne que l'homomorphisme $\widehat A\to\widehat B$ est injectif, et si $k(x)=k(y)$ ou si $k$ est séparablement clos, cet homomorphisme est surjectif par (17.4.4). La réciproque est immédiate.

\begin{proposition}[17.6.4]
\label{IV.17.6.4}
Sous les hypothèses de (17.6.3), si $f$ est étale au point $x$, on a $\dim(\sh{O}_{X,x})=\dim(\sh{O}_{Y,y})$.
\end{proposition}

C'est un cas particulier de (17.5.8, (i)) puisque $x$ est isolé dans sa fibre $f^{-1}(y)$.

